% GENERATED FILE. Edit canonical lessons, then run npm run generate:books.
% canonical-source-hash: fnv1a64:47c08dc118b5edb2
% canonical-lessons: ZH-C344-baocun, ZH-C344-jiyi, ZH-C344-wangshi, ZH-C344-xiangce, ZH-C344-yongyuan

\chapter{Keeping Memories}
\label{ch:zh-keeping-memories}
\hlchaptermodality{344}

\begin{chapteropening}
\textbf{By the end of this chapter:} \emph{I can say I kept something, talk about memories, the past and a photo album, and say I will never forget.}

Complete the last lesson of chapter 344: I can say I kept something, talk about memories, the past and a photo album, and say I will never forget.
\end{chapteropening}

\section[bǎocún]{bǎocún — to keep; to save}

\label{lesson:ZH-C344-baocun}



\begin{quote}
Before the new one: say the Mandarin for good wishes sent to someone, then the Mandarin for to share good things.
\end{quote}

\subsection*{You'll want to know: bǎocún}
\emph{bǎocún} — \textquotedblleft{}to keep; to save\textquotedblright{}.

\begin{quote}
\emph{tā de xìn wǒ dōu bǎocún zài shūfáng lǐ} — I've kept all his letters in the study.
\end{quote}

Its characters come later; for now it is a word for the ear.

\begin{cousinweb}
\emph{bǎocún} is built from \emph{bǎo} and \emph{cún}, both inherited from Middle Chinese, \emph{pawX} and \emph{dzwon} (Guangyun, in Baxter's transcription), which glosses the first 'to keep safe' and the second 'to be there'.
\end{cousinweb}

\begin{grammarlens}[title={What you've built}]
One more word for telling what happened.
\end{grammarlens}

\subsection*{Your turn}
\begin{itemize}
\raggedright
  \item \emph{Say it:} \emph{bǎocún}
  \item \emph{Say it:} \emph{bǎocún}, once more
  \item \emph{Say it:} \emph{tā de xìn wǒ dōu bǎocún zài shūfáng lǐ} — I've kept all his letters in the study.
  \item \emph{Recall:} say the Mandarin for good wishes sent to someone, then the Mandarin for to share good things, then say \emph{bǎocún} again
\end{itemize}

\begin{tcolorbox}[breakable,colback=teal!4,colframe=teal!35!black,title={Before you move on}]
What does bǎocún mean? (\textquotedblleft{}To keep; to save\textquotedblright{}.) Say it once more.
\end{tcolorbox}
\section[jìyì]{jìyì — memory}

\label{lesson:ZH-C344-jiyi}



\begin{quote}
Before the new one: say the Mandarin for to share good things, then the Mandarin for to keep something safe.
\end{quote}

\subsection*{You'll want to know: jìyì}
\emph{jìyì} — \textquotedblleft{}memory\textquotedblright{}.

\begin{quote}
\emph{nà shì wǒ zuì hǎo de jìyì} — That's my best memory.
\end{quote}

Its characters come later; for now it is a word for the ear.

\begin{cousinweb}
\emph{jìyì} is built from \emph{jì} and \emph{yì}, both inherited from Middle Chinese, \emph{kiH} and \emph{'ik} (Guangyun, in Baxter's transcription), which glosses the first 'to record' and the second 'to think of'.
\end{cousinweb}

\begin{grammarlens}[title={What you've built}]
One more word for telling what happened.
\end{grammarlens}

\subsection*{Your turn}
\begin{itemize}
\raggedright
  \item \emph{Say it:} \emph{jìyì}
  \item \emph{Say it:} \emph{jìyì}, once more
  \item \emph{Say it:} \emph{nà shì wǒ zuì hǎo de jìyì} — That's my best memory.
  \item \emph{Recall:} say the Mandarin for to share good things, then the Mandarin for to keep something safe, then say \emph{jìyì} again
\end{itemize}

\begin{tcolorbox}[breakable,colback=teal!4,colframe=teal!35!black,title={Before you move on}]
What does jìyì mean? (\textquotedblleft{}Memory\textquotedblright{}.) Say it once more.
\end{tcolorbox}
\section[wǎngshì]{wǎngshì — past events; things of the past}

\label{lesson:ZH-C344-wangshi}



\begin{quote}
Before the new one: say the Mandarin for to keep something safe, then the Mandarin for a memory of the past.
\end{quote}

\subsection*{You'll want to know: wǎngshì}
\emph{wǎngshì} — \textquotedblleft{}past events; things of the past\textquotedblright{}.

\begin{quote}
\emph{nà tiān wǎnshang yéye hé wǒmen shuōqǐ le wǎngshì} — That evening my grandpa talked to us about the old days.
\end{quote}

Its characters come later; for now it is a word for the ear.

\begin{cousinweb}
\emph{wǎngshì} is built from \emph{wǎng} 'to go' and \zh{事} 'a matter', both inherited from Middle Chinese, \emph{hjwangX} and \emph{dzriH} (Guangyun, in Baxter's transcription), which glosses the first 'to go'.
\end{cousinweb}

\begin{grammarlens}[title={What you've built}]
One more word for telling what happened.
\end{grammarlens}

\subsection*{Your turn}
\begin{itemize}
\raggedright
  \item \emph{Say it:} \emph{wǎngshì}
  \item \emph{Say it:} \emph{wǎngshì}, once more
  \item \emph{Say it:} \emph{nà tiān wǎnshang yéye hé wǒmen shuōqǐ le wǎngshì} — That evening my grandpa talked to us about the old days.
  \item \emph{Recall:} say the Mandarin for to keep something safe, then the Mandarin for a memory of the past, then say \emph{wǎngshì} again
\end{itemize}

\begin{tcolorbox}[breakable,colback=teal!4,colframe=teal!35!black,title={Before you move on}]
What does wǎngshì mean? (\textquotedblleft{}Past events; things of the past\textquotedblright{}.) Say it once more.
\end{tcolorbox}
\section[xiàngcè]{xiàngcè — photo album}

\label{lesson:ZH-C344-xiangce}



\begin{quote}
Before the new one: say the Mandarin for a memory of the past, then the Mandarin for past events.
\end{quote}

\subsection*{You'll want to know: xiàngcè}
\emph{xiàngcè} — \textquotedblleft{}photo album\textquotedblright{}.

\begin{quote}
\emph{wǒ zài xiàngcè lǐ kàndào le xiǎoshíhou de zhàopiàn} — In the album I came across photos of myself as a child.
\end{quote}

Its characters come later; for now it is a word for the ear.

\begin{cousinweb}
\emph{xiàngcè} is built from \emph{xiàng} 'a likeness' and \emph{cè} 'a volume', both inherited from Middle Chinese, \emph{sjangH} and \emph{tsrheak} (Guangyun, in Baxter's transcription), which glosses the first 'to look at' and the second 'bound bamboo slips'. The sense 'a photograph' for the first is modern.
\end{cousinweb}

\begin{grammarlens}[title={What you've built}]
One more word for telling what happened.
\end{grammarlens}

\subsection*{Your turn}
\begin{itemize}
\raggedright
  \item \emph{Say it:} \emph{xiàngcè}
  \item \emph{Say it:} \emph{xiàngcè}, once more
  \item \emph{Say it:} \emph{wǒ zài xiàngcè lǐ kàndào le xiǎoshíhou de zhàopiàn} — In the album I came across photos of myself as a child.
  \item \emph{Recall:} say the Mandarin for a memory of the past, then the Mandarin for past events, then say \emph{xiàngcè} again
\end{itemize}

\begin{tcolorbox}[breakable,colback=teal!4,colframe=teal!35!black,title={Before you move on}]
What does xiàngcè mean? (\textquotedblleft{}Photo album\textquotedblright{}.) Say it once more.
\end{tcolorbox}
\section[yǒngyuǎn]{yǒngyuǎn — forever}

\label{lesson:ZH-C344-yongyuan}



\begin{quote}
Before the new one: say the Mandarin for past events, then the Mandarin for a photo album.
\end{quote}

\subsection*{You'll want to know: yǒngyuǎn}
\emph{yǒngyuǎn} — \textquotedblleft{}forever\textquotedblright{}.

\begin{quote}
\emph{wǒ huì yǒngyuǎn jìde nàge Zhōngqiū} — I'll always remember that moon festival.
\end{quote}

Its characters come later; for now it is a word for the ear.

Two third tones meet in this word: yǒng is said with a rising tone before yuǎn, and the spelling keeps both marks.

\begin{cousinweb}
\emph{yǒngyuǎn} is built from \emph{yǒng} and \emph{yuǎn} 'far', both inherited from Middle Chinese, \emph{hjwaengX} and \emph{hjwonX} (Guangyun, in Baxter's transcription), which glosses the first 'long' and the second 'distant'.
\end{cousinweb}

\begin{grammarlens}[title={What you've built}]
One more word for telling what happened.
\end{grammarlens}

\subsection*{Your turn}
\begin{itemize}
\raggedright
  \item \emph{Say it:} \emph{yǒngyuǎn}
  \item \emph{Say it:} \emph{yǒngyuǎn}, once more
  \item \emph{Say it:} \emph{wǒ huì yǒngyuǎn jìde nàge Zhōngqiū} — I'll always remember that moon festival.
  \item \emph{Recall:} say the Mandarin for past events, then the Mandarin for a photo album, then say \emph{yǒngyuǎn} again
\end{itemize}

\begin{tcolorbox}[breakable,colback=teal!4,colframe=teal!35!black,title={Before you move on}]
What does yǒngyuǎn mean? (\textquotedblleft{}Forever\textquotedblright{}.) Say it once more.
\end{tcolorbox}
